% !TEX root = ../../main.tex

\section{本章小结}

本章围绕 HearBridge 系统核心功能实现展开说明。
首先，介绍了 HarmonyOS 移动端的启动登录、资源浏览、实时训练交互和句子视频识别页面实现；
其次，说明了 Spring Boot 服务端的分层结构、认证登录态维护、资源服务、句子视频识别服务、DeepSeek 语义修正服务、文件上传与资源访问以及关键接口调用；
随后，介绍了 Vue 后台管理端在登录认证与路由控制、分类资源管理、样本与模型训练管理以及模型版本管理方面的实现；
接着，说明了模型版本发布与 Python 识别服务重载之间的联动过程；
最后，分析了系统异常处理与用户反馈机制。

通过本章实现说明，可以看出 HearBridge 系统并非由单一端完成全部功能，而是通过移动端、后台管理端、Spring Boot 服务端和 Python 识别服务的协同实现核心业务闭环。
下一章将在系统实现的基础上，对手势识别服务实现过程和识别实验结果进行进一步分析。
